%%%%%%%%%%%%%%%%%%%%%%%%%%%%%%%%%%%%%%%%%%%%%%%%%%%%
%    Canadian AI Latex Template    %
%%%%%%%%%%%%%%%%%%%%%%%%%%%%%%%%%%%%%%%%%%%%%%%%%%%%
\documentclass[10pt]{cai26}

\begin{document}
% Editorial staff will replace the following values:
% 1. Conference Year
% 2. Issue number
% 3. Article DOI
\def\conferenceyear{2026}
\volumeheader{39}{0}%{00.000}
\begin{center}

\title{Temporal Characterization and Deep Learning-Based Modeling of Human Sleep from Multivariate Polysomnography-Derived Features}
\maketitle

\thispagestyle{empty}
\pagenumbering{gobble}

% Add Authors and Affiliations in the camera ready
% for the double blind review, please leave this section as is 
\begin{tabular}{cc}
First Author\upstairs{\affilone,*}, Second Author\upstairs{\affilone}, Third Author\upstairs{\affilthree}
\\[0.25ex]
{\small \upstairs{\affilone} Affiliation One} \\
{\small \upstairs{\affiltwo} Affiliation Two} \\
{\small \upstairs{\affilthree} Affiliation Three} \\
\end{tabular}
  
% Replace with corresponding author email address
\emails{
  \upstairs{*} corresponding\_author@example.ca 
}
\vspace*{0.1in}
\end{center}

\begin{abstract}
More recently, artificial intelligence is applied in studying physiological signals to evaluate health, yet the modeling of temporal dependencies is still a focus of 
research. In sleep studies, for example, 30 second epochs are frequently treated as independent samples, preventing the study of sleep stage transitions, signal 
interactions, and broader patterns in an individual's sleep architecture. To address these limitations, the present work presents a modular temporal analysis for 
modeling sleep physiology in a model based on features derived from polysomnography (PSG). The method combines subject-wise sleep pattern profiling with multivariate 
temporal dependency modeling and sequence-based sleep stage prediction to get both short and long range temporal dynamics in EEG, EOG, and EMG signals. Linear vector 
autoregression is used to quantify cross signal relationships and temporal persistence, while a Bi-LSTM sequence model capitalizes on contextual information for sleep 
stage classification. Experiments using the Sleep-EDF dataset show substantial differences across individuals in sleep fragmentation, strong within-signal temporal 
continuity, and selective cross-modal coupling. Adding temporal structure improves classification accuracy from 74.92\% to 83.94\% and macro F1-score boosts up from 
0.57 to 0.65 due to better detection of minority and transitional stages. For all these reasons, the results highlight an increase in interpretability along with 
improved predictive performance of temporal modeling, further supporting its use for enhanced sleep-stage characterization and structured physiological analysis.
\end{abstract}

% add your keywords
\begin{keywords}{Keywords:}
Sleep analysis, Polysomnography, Temporal modeling, Sleep-stage transitions, Vector autoregression, Recurrent neural networks, EEG, EOG, EMG
\end{keywords}
\copyrightnotice

%{\noindent{\bf Editors:} Lydia Bouzar-Benlabiod, Carson K. Leung}

\section{Introduction}
\label{introduction}
Sleep is a structured neurophysiological activity which is heavily determined by the interplay of changes in the neural, ocular, and muscular systems. Polysomnography 
(PSG) is a type of monitoring that provides synchronized measurements from electroencephalography (EEG), electrooculography (EOG), and electromyography (EMG), thus 
enabling detailed tracking of the sleep dynamics over the duration of 30 second epochs~\cite{Kemp2000SleepEDF}. Cumulatively these signals provide representation of the progression of sleep 
stages, transitions, and global sleep organization throughout the night.

Recent developments in deep learning systems can learn representations directly from physiological signals, greatly enhancing automatic sleep-stage classification~\cite{Supratak2017,Chambon2018}. Many 
existing approaches focus on per-epoch accuracy but consider stage labels as independent objectives~\cite{Supratak2017}. These formulations frequently ignore the structural temporal 
dependencies that characterize the evolution of sleep such as probability-based stage transitions, continuity conditions, and cyclic architecture structures~\cite{Phan2019}. Modeling 
sleep as a temporally coherent multivariate process continues to be challenging.

Apart from the accuracy of classification, insights into sleep-stage transition behavior and global sleep architecture are also important for accounting for 
inter-subject variability and recognizing stable sleep phenotypes~\cite{Phan2019J,Phan2022}. Physiological signals captured during sleep exhibit a dynamic temporal structure that is amenable not 
only to predictive modeling but also statistical interpretation of dynamic patterns dependent on modalities. Combining sequential deep learning with interpretable 
temporal analysis paves the way for a comprehensive modeling of sleep organization.

This study proposes a temporally relevant modeling scheme to mitigate these limitations of PSG-based physiological signatures. Rather than treating sleep stages as 
independent classification targets, the proposed approach views sleep as a structured evolving multivariate sequence influenced by cross-modal physiological interactions 
and transition dynamics~\cite{Chambon2018,Phan2022}. It combines sequence-based predictive modeling with statistical temporal analysis in order to characterize dependencies among EEG, EOG, and EMG 
data, and stage transition behavior, as well as to analyze global sleep architecture patterns across the subjects.

This work is useful in three ways. First, we describe subject-level sleep architecture through performance assessments across different individuals, summarizing and 
comparing structure, continuity, and fragmentation of stages in human sleep. Second, we build a sequence-based model of sleep-stage prediction from an epoch-wise 
perspective that allows taking into account temporal dependencies in multivariate physiological signals~\cite{Supratak2017,Phan2019,Pham2023}. Third, we investigate time dependence between data signals and 
potential stage transitions to help interpret the temporal evolution of sleep.

By incorporating linear statistical modeling and nonlinear sequence-based learning, this work progresses the analysis of sleep physiology based on the epoch-level 
classification beyond independent epoch-level classification. This approach focuses on reproducibility and scalability, allowing for systematic evaluation of temporal 
structure in PSG-generated signals.

The proposed methodology operates at three complementary temporal resolutions. We use subject-wise aggregation to characterize inter-individual variability in sleep 
organization at the population level. At the within-subject level, temporal dependency modeling models short-term physiological interactions in sleep. At the epoch level, 
sequence-based predictive modeling considers how contextual information affects sleep-stage classification. This multi-scale design allows for concurrent examination of 
sleep dynamics at population, individual, and epoch-level levels.

\section{Literature Review}
\label{literature}
Automated sleep-stage classification has also evolved, driven by deep learning architectures that learn hierarchical and temporal 
representations directly from polysomnography (PSG) signals. Initial computational techniques typically based on time-domain and 
spectral features generated manually from EEG, EOG, and EMG signals, along with classical classifiers, including support vector 
machines and ensemble models. Although capable of implementing these methods, 30 second epochs were considered as independent 
observations and did not reflect the structured temporal evolution in sleep architecture. 

With the advent of deep neural networks, this paradigm changed into a sequence-aware modeling. DeepSleepNet~\cite{Supratak2017} showed that 
integration of convolutional neural networks (CNNs) and bidirectional long short-term memory (Bi-LSTM) networks facilitated the 
automatic extraction of features and the modeling of contextual temporal relationships. While the local characteristics of signal 
are encoded by the convolutional element, the complex inter-epoch dependencies via the recurrent element minimize the possibility 
of introducing unrealistic stage transitions, and therefore enhancing the stability of the system. Expanding on this concept a 
hierarchical sequence-to-sequence architecture of a multi-dimensional architecture was introduced in SeqSleepNet~\cite{Phan2019}, for 
jointly modeling intra-epoch representations and inter-epoch transitions. As an interpretation, the model further enhanced sleep 
staging as a sequential prediction problem, because they predicted non-isolated stage sequences instead of separate labels from 
each other. 

Later studies explored hybrid temporal architectures. In this manner, CNN-Transformer-LSTM models~\cite{Pham2023} combined attention 
mechanisms to capture long-range dependencies and convolutional as well as recurrent processing. Joint 
classification-and-prediction schemes~\cite{Phan2019J} added adjacent prediction epochs in order to achieve local transition smoothness 
without large input sequences. Combined ensemble methods that employ many end-to-end sequence models achieved enhanced robustness 
and associated transition sensitivity~\cite{Wang2020}. DSSNet~\cite{Chang2022} focus was on deep sequential representation learning, so as to 
highlight the need to consider temporal structure beyond straightforward epoch classification. 

Integrated multimodal and multi-view modeling approaches that exploit complementary physiological signals have been reinforced. 
Prior studies have proposed context-based multimodal CNN-BiLSTM architectures~\cite{Chambon2018} and multi-channel CNN-BiLSTM hybrid 
networks~\cite{Karim2025} and their combined model for EEG, EOG, and EMG signals improved discrimination of REM, non-REM, and wake state, 
especially in transitional borderline states. XSleepNet~\cite{Phan2022} proposed a multi-view sequential framework that unites raw 
waveform signals with time-frequency representations, performing adaptive gradient blending to adjust heterogeneous feature 
spaces. In CareSleepNet~\cite{Wang2024}, convolutional and attention-based approach were further integrated to accelerate cross-modal 
fusion and temporal sensitivity.

Taken together, these developments provide three notable trends in contemporary automated analysis of sleep: (i) end-to-end 
representation learning, (ii) modeling explicitly time sequence, and (iii) multimodal physiology integration. These architectures 
show the impact of adding temporal context on predictive performance and transition continuity. 

Despite significant advancements in sequence modeling and multimodal fusion have been made, many current techniques focus on 
providing classification accuracy metrics only, failing to embed temporal dependencies implicitly in deep networks. Explicit 
statistical measurements of cross modal physiological interplay and lag based temporal persistence are still relatively scarce. 
Moreover, feature-based characterization of sleep architecture subject-level features including fragmentation patterns, continuity 
measures, and inter individual variability are rarely considered when implementing predictive prediction processes. Deep sequence 
models capture contextual information, but do not generally provide interpretable estimates of cross-signal interactions, and 
structured temporal persistence.

Consequently, previous studies have largely contributed to advancing sequential modeling~\cite{Supratak2017,Phan2019}, 
hybrid models~\cite{Pham2023,Phan2019J,Wang2020,Chang2022} and multimodal integration~\cite{Chambon2018,Karim2025,Phan2022,Wang2024}, 
an integrated model which connects interpretable 
multivariate temporal modeling to deep sequence learning and subject-level sleep architecture characterization, a unified 
pipeline linking interpretable multivariate temporal modelling and deep sequence learning and subject-level sleep architecture 
characterization is still largely underinvestigated. Closing this gap allows for the modeling of sleep as not only a discrete 
classification task, but also an organized, multi-scaled and coordinated physiological process governed by neural, ocular, and 
muscular dynamics.

\section{Methodology}
\label{method}
The analysis is structured around three complementary components: (i) subject-wise sleep pattern analysis, (ii) within-subject temporal dependency and transition 
modeling, and (iii) epoch-level predictive sequence modeling.

\subsection{Dataset and Data Representation}

Analyses were performed using the Sleep-EDF Expanded database~\cite{Kemp2000SleepEDF,Goldberger2000PhysioNet}, a publicly available benchmark dataset widely used in automatic sleep staging research. The 
dataset contains full-night PSG recordings collected from adult subjects under laboratory conditions and annotated according to standard clinical sleep-scoring 
guidelines~\cite{AASM2014}.

Each recording consists of synchronized EEG, EOG, and EMG signals segmented into non-overlapping 30-second epochs. EEG signals are obtained from Fpz-Cz and Pz-O
z derivations, EOG from the horizontal channel, and EMG from the submental channel. For each epoch and channel, statistical features (mean, standard deviation, RMS) are 
extracted, yielding multivariate time series indexed by epoch. Sleep stages are annotated at the epoch level with five classes: Wake (W), N1, N2, N3, and REM.

All physiological features are standardized per subject to reduce inter-individual variability. EMG features are further normalized using a rolling-window z-score to 
address amplitude variability and transient muscle activity. Temporal ordering is preserved throughout all analyses.

\subsection{Subject-Wise Sleep Pattern Analysis}

Aggregated temporal representations summarise and compare sleep architecture across individuals. Epoch level sleep stage sequences 
are summarized on the recording and subject levels to capture inter subject variability in sleep organization. For every night 
recording of the subject, sleep fragmentation is measured across two metrics~\cite{Bonnet2003Fragmentation}; sleep stage transitions (the change from consecutive 
epoch labels), and awakenings (the sleep stage transition from any sleep stage N1, N2, N3, or REM to Wake). These metrics are 
calculated directly from the ordered sleep stage sequence and averaged to derive subject level means over nights.

Global sleep organization is further described by mean metrics using summary measures, which derive from epoch-level labels. Sleep 
efficiency (SE\%) is computed as the proportion of epochs classified as sleep relative to the total number of epochs in a 
recording~\cite{Bonnet2003Fragmentation}. Deep sleep proportion (N3\%) is computed as the percentage of N3 epochs relative to the total number of sleep epochs. 
Collectively, these normalized metrics serve to make the sleep architecture at the subject level more clearly apparent for 
subsequent comparisons between subjects. The subject-level fragmentation measurements are also used to identify representative 
recordings for fine temporal modeling, where individuals with the highest fragmentation are given specific attention.

\subsection{Temporal Dependency and Transition Analysis}

To examine short-term physiological interactions and stage transition dynamics~\cite{Taguchi2000Markov}, Vector Autoregression (VAR)~\cite{Lutkepohl2005VAR} is employed to jointly model EEG, EOG, and EMG dynamics.

\begin{equation}
\mathbf{y}_t =
\begin{bmatrix}
EEG_t \\
EOG_t \\
EMG_t
\end{bmatrix}
\end{equation}

A VAR($p$) model is defined as:

\begin{equation}
\mathbf{y}_t = \mathbf{c} + \sum_{k=1}^{p} \mathbf{A}_k \mathbf{y}_{t-k} + \boldsymbol{\varepsilon}_t
\end{equation}

Stationarity assumptions are evaluated using the Augmented Dickey-Fuller (ADF) test~\cite{DickeyFuller1979ADF} and the KPSS test~\cite{Kwiatkowski1992KPSS}. Model parameters are estimated via ordinary 
least squares~\cite{Lutkepohl2005VAR}.

\subsection{Predictive Sequence Modeling}

Sleep staging is restructured and represented as a temporal sequence modeling problem, in which consecutive epochs are grouped 
into windows of length $T = 20$, corresponding to a 10 minute temporal context. The model predicts the sleep stage label 
associated with the final epoch in the sequence given an ordered sequence of physiological feature vectors 
${\mathbf{x}_1, \mathbf{x}_2, \ldots, \mathbf{x}_T}$. This formulation enables the classifier to learn contextual 
information from earlier epochs rather than treating each epoch independently. 

All input features are standardized using statistics only calculated from the training data to avoid leakage of information and 
to ensure uniform scaling of subjects. Temporal dependencies are modeled with a two layer bidirectional Long Short Term Memory 
network~\cite{Hochreiter1997LSTM,Schuster1997BiLSTM}, that encompasses nonlinear relationships in both forward and backward temporal directions. Then, the resulting 
hidden representation is passed through a fully connected classification layer, and class probabilities are obtained using a 
Softmax activation function~\cite{Goodfellow2016DeepLearning}. 

To compensate for class imbalance through sleep stages, training makes use of a class weighted cross entropy loss~\cite{Goodfellow2016DeepLearning}, in 
which higher weight gets assigned to the underrepresented stages. This optimization is done by AdamW~\cite{Loshchilov2019AdamW} with a weight 
decay of $1 \times 10^{-4}$, and dropout is done with a rate of 0.3 between recurrent layers to prevent overfitting, while also 
enabling better generalization as well.

\section{Results and Discussion}
\label{results}

This section presents the findings related to subject wise sleep pattern analysis, temporal dependency modeling, and sequence based sleep stage prediction.

\subsection{Subject-Wise Sleep Pattern Analysis}

Sleep fragmentation was evaluated across all subjects using two complementary quantitative measures. The first measure is the average number of sleep stage transitions 
per night, computed as the number of changes between consecutive epoch level stage labels. The second measure is the average number of awakenings, defined as transitions 
from any sleep stage including N1, N2, N3, or REM to Wake. Together, these metrics provide a structured assessment of sleep continuity and stage instability at the 
subject level.

\begin{figure}[H]
\centering
\includegraphics[width=0.85\linewidth]{figs/fragmentation_top15.png}
\caption{Top 15 subjects ranked by average number of sleep-stage transitions, illustrating inter-subject variability in sleep fragmentation.}
\label{fig:fragmentation}
\end{figure}

Figure 1 presents the top 15 most fragmented subjects based upon the average number of stage transitions. A lot of inter-subject 
variability occurs, and fragmentation scales from around 200 to over 350 transitions per night. The highest level of fragmentation 
was recorded in Subject 4652, implying a very unstable sleep architecture. 

In a concurrent study, the association was analyzed of fragmentation and sleep efficiency. Across respondents, there was a 
moderate level of a positive relationship between number of awakenings and sleep efficiency 
(Pearson $r = 0.301$, Spearman $\rho = 0.297$). Though increased fragmentation is typically linked to decreased sleep quality, 
this finding indicates considerable heterogeneity in terms of the compensating mechanisms that individuals use to make up for 
frequent awakenings and reiterates the importance of subject-level analysis rather than population averages. 

The second night recording of the Subject 465 represents an extreme fragmentation phenotype which was selected for detailed '
temporal modeling in later analyses based on these findings.

\subsection{Temporal Dependency and Transition Analysis (VAR)}

A Vector Autoregression (VAR) model was fitted to EEG Fpz-Cz RMS, EEG Pz-Oz RMS, EOG horizontal RMS, and EMG submental RMS signals 
for subject 465, night 2. Four lags corresponding to 30 seconds, 1 minute, 2 minutes, and 4 minutes were included to capture 
short-term physiological interactions. 

Strong within-signal temporal persistence was demonstrated, especially in regard to EMG recordings. Past EMG values strongly 
predicted current EMG values (lag 1 coefficient = 0.715, $p < 0.001$), suggesting high short-term muscle tone continuity during 
sleep. EEG signals also showed strong autoregressive structure, reflecting rhythmic neural activity in sleep stages. 

Specific cross-modal interactions were evident between groups. Frontal EEG activity somewhat predicted later EOG horizontal 
activity (lag 1 coefficient = 0.042, $p = 0.028$). This indicates short-term coupling between neural and ocular activity. In 
contrast, cross-modal influences involving EMG were weaker, indicating the muscles were quite autonomous during periods of 
sleep. 

Residual correlations were moderate between EEG and EOG ($r \approx 0.45$) and low between EMG and other modalities 
($r < 0.16$), suggesting that the VAR model captured a great deal of shared temporal structure. 

Table 1 summarizes forecast accuracy for each physiological signal. EEG signals have the lowest prediction error overall, and 
posterior EEG (Pz-Oz) shows the highest level of short-term predictability. Similar results for EMG submental activity also 
indicate low forecast error indicative of high temporal persistence of muscle tone during sleep.

\begin{table}[H]
\centering
\begin{tabular}{lcc}
\toprule
\textbf{Signal} & \textbf{MAE} & \textbf{RMSE} \\
\midrule
EEG Fpz--Cz (RMS) & 1.362 & 1.678 \\
EEG Pz--Oz (RMS) & 0.536 & 0.640 \\
EOG Horizontal (RMS) & 1.528 & 1.729 \\
EMG Submental (RMS, normalized) & 0.472 & 0.619 \\
\bottomrule
\end{tabular}
\caption{Forecast accuracy of VAR model for subject 465, night 2. MAE: Mean Absolute Error; RMSE: Root Mean Squared Error.}
\label{tab:var_accuracy}
\end{table}

Horizontal EOG, in comparison, has much larger errors, which is indicative of its rapid fluctuation and sensitivity to 
sleep-stage transitions. These differences suggest that short-term predictability differs greatly by modality, with EEG and 
EMG following more stable timelines than ocular activity.

Figure 2 compares observed and VAR-forecasted physiological signals for subject 465, night 2. EEG and EMG dynamics are well 
observed by the model but larger deviation from the predicted behavior is seen for horizontal EOG activity, consistent with its 
irregular changes. 

\begin{figure}[H]
  \centering
  \includegraphics[width=\columnwidth]{figs/forecasting.png}
  \caption{Observed and VAR-forecasted EEG, EOG, and EMG RMS signals for subject 465, night 2.}
  \label{fig:var_forecast}
\end{figure}

Overall, the results confirm that physiological sleep dynamics are signal-specific and are structured 
rather than stochastic, and indicate the advantages of multivariate temporal models for sleep analysis.

\subsection{Predictive Sequence Modeling (Bi-LSTM)}

The training and validation behavior of the proposed Bi-LSTM model is shown in Fig. 3. Training loss decreases steadily along 
consecutive epochs, while validation accuracy stabilizes, suggesting convergence and generalization to unseen subjects 
(with class-weighted tuning). 

\begin{figure}[H]
  \centering
  \includegraphics[width=\columnwidth]{figs/training_validation.png}
  \caption{Bi-LSTM temporal learning process.}
  \label{fig:bilstm_train}
\end{figure}

While validation accuracy is stable, the validation loss is often greater than the training loss. This highlights the striking 
class imbalance in sleep-stage distributions where Wake (W) is predominant while transitional stages like N1 and REM occur less 
frequently than the normal stages. Class-weighted optimization raises the penalty for misclassifications of minority stages; as 
such, a relatively small number of errors in these stages can yield a relatively large validation loss even if overall accuracy 
remains stable. 

Furthermore, validation recordings refer to previously unseen subjects, creating physiological variability in the EEG, EOG, and 
EMG signal characteristics. The inter-subject variability further adds uncertainty around sleep-stage transitions, where 
classification mistakes are both more frequent and more severely penalized. Consequently, elevated validation loss indicates 
stronger sensitivity toward minority and transitional stages rather than reduced model generalization. 

In Table 2, the quantitative performance on the validation set is compared between the sequence-based Bi-LSTM model and the 
epoch-level XGBoost baseline.

\begin{table}[H]
\centering
\caption{Performance comparison of sleep-stage classification models.}
\label{tab:model_perf}
\renewcommand{\arraystretch}{1.15}
\setlength{\tabcolsep}{6pt}
\begin{tabular}{|l|c|c|c|c|}
\hline
\textbf{Model} & \textbf{Accuracy} & \textbf{Prediction} & \textbf{Recall} & \textbf{F1 Score} \\
\hline
XGBoost (epoch-level) & 74.92\% & 0.5590 & 0.6225 & 0.5745 \\
Bi-LSTM (sequence-level) & \textbf{83.94\%} & \textbf{0.6721} & \textbf{0.6325} & \textbf{0.6474} \\
\hline
\end{tabular}
\end{table}

The Bi-LSTM shows an absolute accuracy gain of 9.02 percentage points over the non-temporal baseline. Macro-averaged precision 
and F1 score improvements suggest increased detection of minority and transitional sleep stages, validating the utility of 
incorporating temporal context into sleep-stage prediction. 

After training, the Bi-LSTM also performs complete automatic sleep-stage labeling at a standard 30 second epoch resolution. During 
inference, sequences of physiological characteristics collected from EEG, EOG, and EMG signals are processed to predict the sleep 
stage of the most recent epoch. The model generates a result by sliding the temporal window over an entire recording, producing a 
fully automatically labeled hypnogram which has direct comparisons to expert-scored annotations. 

In addition to discrete sleep-stage predictions made, the model also provides a confidence score for each epoch, defined as the 
maximum predicted class probability. In stable sleep conditions, higher confidence values are usually seen, while lower confidence 
values occur most often near sleep-stage transitions and physiologically ambiguous stages such as N1 and REM. This measure of 
confidence gives a readable indication of the reliability of prediction and emphasizes epochs with increased temporal uncertainty.

\subsection{Discussion}

These findings confirm that sleep physiology is regulated by organized temporal dependencies that are multi-level effects such as 
subject-wise structure, inter-signal dynamics, and sequence-level prediction~\cite{Phan2019,Phan2022,Li2024SleepFC}.

The data shows significant heterogeneity in sleep fragmentation that confirms population-averaged metrics are not ideal for study. 
Highly fragmented individuals allow for targeted temporal analysis and assist in personalized sleep characterization.

VAR analysis verifies that EEG, EOG, and EMG signals display strong within-signal persistence and selective cross-modal coupling, 
as demonstrated by VAR modeling. Neural-ocular interactions can take place over short time spans, and muscle activity is largely 
self-driven. These results give direct empirical support that sleep-stage transitions arise from coordinated, yet asymmetric 
physiological interactions rather than independent processes~\cite{Chen2019LFP}.

Modeling by sequence also shows that embedding temporal context substantially enhances sleep-stage classification~\cite{Supratak2017,Mousavi2019SleepEEGNet}. The superiority 
of the Bi-LSTM over epoch-based baselines indicates that sleep needs to be considered as a continuous process rather than just 
decisions made in isolation. Increased validation loss with stable accuracy indicates sensitivity to minority stages not 
decreased generalization.

The results collectively provide a coherent picture of sleep as a multivariate, time-dependent system in which subject-wise 
structure, short-term physiological dynamics, and long-range temporal context combine to determine the patterned sleep 
architecture.

\section{Conclusion}

The work presented a unified framework to assess human sleep as a multivariate and temporally organized physiological process 
with cues derived from polysomnography. By going beyond independent epoch-level analysis, the proposed approach identifies sleep 
dynamics at the basis of (1) subject-wise architecture (2) inter-signal temporal dependencies (3) sequence-based sleep-stage 
prediction. Results thus establish that sleep physiology presents an interpretable temporal structure throughout sleep, consistent 
with the view of sleep as continuous and regulated by coordinated neural, ocular, and muscular activity. 

Subject-wise analysis showed that sleep fragmentation varies greatly between individuals, which suggests that population averages 
do not fully capture personal sleep patterns. Identifying highly fragmented subjects allowed for more focused temporal modeling and 
helped in selecting representative cases for detailed study. These findings highlight the importance of analyzing sleep at the 
individual level to better understand sleep dynamics.

Temporal dependency analysis by vector autoregression corroborated considerable within signal persistence as well as selective 
cross-modal encounters of EEG, EOG, and EMG signals. Thus, the mismatch in cross-signal coupling we observed indicates that 
sleep-stage transitions result from coordinated but uneven physiological interactions, rather than independent or purely top-down 
processes. Forecasting also proved that short-term physiological dynamics are well established signal-specific and predictable, 
which further supports the applicability of time profiling for sleep analysis. 

Bi-LSTM-based sequential modeling significantly enhanced sleep-stage classification versus epoch-level baselines. By factoring in 
temporal context the model attained superior accuracy and recognition of minority and transient sleep phases. Predictions were 
even more confidence-aware to enable interpretable markers of time uncertainty, especially in the vicinity of sleep-stage borders. 

Collectively, these results show that using temporal sequence modeling leads to higher and more stable sleep-stage prediction compared to 
epoch-level methods. This study combines subject wise analysis, interpretable temporal modeling, and sequence-based prediction in a 
single approach. The findings highlight the importance of reproducibility and clear physiological interpretation in multivariate sleep 
analysis. Future work may include adding more physiological signals, studying longer time periods, and developing scalable methods for 
personalized sleep monitoring in real-world settings.

\section*{Acknowledgements}
This section can be left blank during double-blind review. 

%Appendixes go here
\appendix

% All references should be stored in the file "references.bib".
% That call to use that file is in "cai.cls". 
% Please do not modify anything below this line.
\printbibliography[heading=subbibintoc]

\end{document}
